\chapter{Conclusion}\label{sec:conclusion}

\begin{itemize}
  \item separation through Director-Actor architecture provides robust, computationally tractable mechanism for maintaining narrative control under user interruption
  \item External temporal monitoring, implemented outside LLM, provides  approach to time-budget enforcement
  \item Pairwise AB/BA evaluation with order-bias correction reveals Director Agent significantly outperforms single-agent Baseline on Narrative Control and Interruption Handling 
  
\end{itemize}
 
 


\section{Limitations}\label{sec:limitations}

Our setup has several limits.
\begin{itemize}
  \item We test only two characters, Olaf and Rocket Raccoon, and mostly on Disney stories, so we cannot say how far the results carry over to other characters and genres.
  \item All variants use the same, small language model, GPT-4o-mini. Larger models may be less prone to drift.
  \item The automated evaluation relies on a single judge model for most comparisons, and scripted user inputs cannot react to what the character says.
  \item The beat generator and the lenient Director have not been tested thoroughly.
  \item The user study has a small sample, only 14 participants who completed it, and limited time per participant.
\end{itemize}

\section{Future Work}

\begin{itemize}
  \item integrate with audio and animation pipelines
  \item make it run on real-time maybe for the director (make it run faster)
  \item introduce memory for switching between stories and characters
  \item test with different types of stories 
  \item further test the beat representation and the lenient director mode
  \item compare directly against NarrativeGenie \cite{kumaran2024narrativegenie}, the closest beat-based system; no public implementation was available, so this would require a reimplementation
  
\end{itemize}